% Flavor catalog per-process template.
% process_id: E004
% family: edm_neutrino
% owner_agent_id: PKA-E004
% writer_agent_id: null
% checker_agent_id: null
% opus_signoff_id: null
% Canonical metadata sidecar: E004.yaml
% Status is controlled by the YAML sidecar status_history, not this comment.

\section{E004: Neutron electric dipole moment}

\subsection*{Process}
\textbf{Standard notation:} \(d_n\).  This entry covers the permanent electric
dipole moment of the neutron, cataloged in the \texttt{edm\_neutrino} family
as a hadronic EDM and CP-violation benchmark.  The observable is flavor
diagonal, but it probes CP-odd quark EDMs, quark chromo-EDMs, the QCD
\(\theta\) term, and the Weinberg three-gluon operator.

\subsection*{PDG / equivalent value(s)}
PDG Live 2026 (\texttt{PDG2026:NeutronEDM}) lists the current direct neutron
EDM limit as
\[
  |d_n| < 1.8\times 10^{-26}\ e\,\mathrm{cm}
  \qquad (90\%~\mathrm{CL}),
\]
from Abel et al. 2020 using ultracold neutrons and Ramsey separated oscillatory
fields at PSI (\texttt{PDG2026:NeutronEDM}).  The PDG table records this as
\(<0.18\) in units of \(10^{-25}\ e\,\mathrm{cm}\); see
\texttt{flavor\_catalog/references/E004/pdg2026\_neutron\_edm\_datablock.txt}.
The associated experimental result is
\[
  d_n = (0.0\pm1.1_{\rm stat}\pm0.2_{\rm sys})\times10^{-26}\ e\,\mathrm{cm},
\]
recorded both by PDG and in the arXiv snapshot
\texttt{flavor\_catalog/references/E004/abel2020\_arxiv2001\_11966.txt}
(\texttt{Abel:NeutronEDM2020}).  PDG also lists the superseded Pendlebury
et al. 2015 direct limit \(|d_n|<3.0\times10^{-26}\ e\,\mathrm{cm}\) at
\(90\%\) CL (\texttt{PDG2026:NeutronEDM}).

\subsection*{Relevance to the RS / anarchic-flavor pipeline}
In anarchic warped or composite-Higgs completions, new complex Yukawa spurions
and KK fermion/gauge loops can generate flavor-diagonal CP-odd dipoles even
when tree-level \(\Delta F=2\) observables are controlled.  The neutron EDM is
therefore a standard hadronic CP benchmark for quark EDMs and chromo-EDMs, and
for threshold-induced Weinberg-operator contributions.  It does not map onto
the existing neutral-meson mixing basis; it would test a different CP-odd
operator lane.

\subsection*{Post-2008 developments}
The 2008 CFW reference point (\texttt{CsakiFalkowskiWeiler:RSFlavor2008})
focused on warped-flavor bounds such as KK-gluon constraints.  Since then, the
direct neutron EDM experimental limit improved from the PDG-listed
Pendlebury-era \(3.0\times10^{-26}\ e\,\mathrm{cm}\) bound to the Abel et al.
2020 \(1.8\times10^{-26}\ e\,\mathrm{cm}\) bound
(\texttt{PDG2026:NeutronEDM}; \texttt{Abel:NeutronEDM2020}).  On the model side,
Koenig, Neubert, and Straub 2014
(\texttt{KoenigNeubertStraub:DipoleComposite2014}) made quark electromagnetic
and chromomagnetic dipoles explicit in composite-Higgs and warped
extra-dimensional language, including the role of ``wrong-chirality'' Yukawas
and the neutron EDM.  On the hadronic side, lattice work such as Bhattacharya
et al. 2022 (\texttt{Bhattacharya:NeutronEDMLattice2022}) now directly targets
the neutron EDM induced by the QCD topological term, the Weinberg three-gluon
operator, and the isovector quark chromo-EDM.

\subsection*{Constraint validity and model dependence}
The experimental limit is a robust null bound on \(d_n\), but translating it
into RS parameters is hadronic- and basis-dependent.  A live interpretation
must specify the UV matching to quark EDM, quark chromo-EDM, \(\theta\), and
Weinberg coefficients; their QCD running and threshold matching; and the
low-energy neutron matrix elements.  This is an EDM-adjacent loop constraint,
not a direct \(\Delta F=2\) or lepton-universality observable.

\subsection*{Code coverage in this repo}
\textbf{NO.}  The required catalog greps and a focused search for
\texttt{neutron[ -]?EDM}, \texttt{nEDM}, \texttt{d\_n}, \texttt{chromo-EDM},
\texttt{CEDM}, \texttt{Weinberg}, and \texttt{three-gluon} over
\texttt{quarkConstraints/}, \texttt{qcd/}, \texttt{flavorConstraints/},
\texttt{neutrinos/}, \texttt{yukawa/}, \texttt{warpConfig/},
\texttt{solvers/}, \texttt{scanParams/}, and \texttt{tests/} found no neutron
EDM implementation.  The only nearby dipole code is the off-diagonal
\(\mu\to e\gamma\) helper in \texttt{flavorConstraints/muToEGamma.py:3},
\texttt{:21}, and \texttt{:81}, which is not a hadronic EDM calculation.

\subsection*{Implementation difficulty}
\textbf{HIGH.}  A production constraint needs new CP-odd dipole and gluonic
operator matching, QCD running with thresholds, and hadronic neutron matrix
elements.  The existing SLL/SLR/VLL/VRR/LR \(\Delta F=2\) infrastructure does
not cover quark EDMs, chromo-EDMs, or the Weinberg operator.

\subsection*{Key references}
Process-local source keys before bibliography consolidation:
\texttt{PDG2026:NeutronEDM}, \texttt{Abel:NeutronEDM2020},
\texttt{CsakiFalkowskiWeiler:RSFlavor2008},
\texttt{KoenigNeubertStraub:DipoleComposite2014}, and
\texttt{Bhattacharya:NeutronEDMLattice2022}.
